\chapter{\babTiga}
\label{bab:3}

In this chapter, we will explain the research methodology we employ, including the big picture of the experiment we conduct, the dataset and knowledge bases, as well as the evaluation metrics.

\section{Research Methodology}
The research methodology we employ is the development and application studies. Specifically, we are using the design science research (DSR) studies paradigm \citep{vomBrocke2020}. This paradigm aims to enhance human knowledge by developing innovative artifacts and generating design knowledge through creative solutions to real-world problems \citep{hevner2004design}. It comprises six main activities: problem identification and motivation, definition of the objectives for a solution, design and development, demonstration, evaluation, and communication. The following is the explanation of each activity, in relation to our work.
\begin{enumerate}
    \item \textbf{Problem identification and motivation}: While the previous GraphRAG framework provided a foundation for leveraging open-source components with knowledge graphs, several limitations were identified that hindered optimal performance and scalability. First, the original system's entity and property extraction methods were limited in their ability to handle diverse and complex question types effectively, particularly for multi-hop reasoning scenarios. Second, the original pipeline architecture lacked runtime-configurable components, expanded fallback mechanisms, and comprehensive semantic logging capabilities needed for adaptability, fault tolerance, and operational transparency when dealing with heterogeneous knowledge graph structures. Third, the absence of specialized model training meant the system could not learn domain-specific patterns for Natural Language to SPARQL (NL2SPARQL) generation, relying solely on few-shot prompting approaches. Finally, dataset generation was constrained in scope and variety, limiting the system's ability to handle diverse query types and adapt to different knowledge domains effectively.
    \item \textbf{Definition of the objectives for a solution}: Building upon the original GraphRAG framework, our objective is to develop a significantly enhanced system that addresses the identified limitations through three main improvements: (1) developing comprehensive dataset generation methods using both template-based and random-walk approaches to create diverse, high-quality training data, (2) implementing specialized model finetuning using QLoRA techniques for NL2SPARQL generation tasks to improve accuracy and reasoning capabilities, and (3) enhancing the original pipeline architecture with a LangGraph-based implementation that provides runtime configurability, expanded fallback mechanisms including Google Search integration, and comprehensive semantic logging while significantly improving adaptability, fault tolerance, and operational transparency through process visualization, RDF-based event tracking, and reviewable system diagrams that enable users to verify and audit the decision-making process.
    \item \textbf{Design and development}: The enhanced GraphRAG architecture is developed as three integrated components using only open-source technologies. The first component implements dual dataset generation approaches: template-based generation for linguistically diverse questions and random-walk generation for comprehensive graph coverage through automatic pattern discovery from RDF data. These datasets are designed to be interchangeable based on user needs and knowledge domain requirements. The second component involves finetuning Large Language Models using QLoRA (Quantized Low-Rank Adaptation) techniques on the generated NL2SPARQL (Natural Language to SPARQL) datasets to improve query generation capabilities and reasoning performance. The third component designs a multi-agent system using LangGraph \footnote{\url{https://langchain.com/langgraph/}} that orchestrates specialized agents for translation, entity and property extraction, strategy selection, verbalization, entity and property search, SPARQL generation, and answer synthesis, with intelligent routing between different processing paths based on query complexity and user settings.
    \item \textbf{Demonstration}: After developing the enhanced system, its functionality is demonstrated through comprehensive testing across multiple scenarios. The finetuned models are evaluated on SPARQL generation tasks using both the baseline QALD-9-Plus dataset and our newly generated synthetic datasets from template-based and random-walk approaches. The dataset generation components are tested for quality, diversity, and coverage across different knowledge domains including Wikidata and local knowledge graphs. The multi-agent architecture is demonstrated through real-time question-answering sessions with process visualization, showing adaptive routing between verbalization and SPARQL generation paths. The complete system is applied to multiple knowledge bases including Wikidata, Curriculum Knowledge Graph, Legal Knowledge Graph, and GESIS Knowledge Graph to enable comprehensive performance evaluation across different domains.
    \item \textbf{Evaluation}: The enhanced system undergoes comprehensive end-to-end quantitative evaluation. System performance is assessed using Jaccard similarity, precision, recall, F1-score, and success rate metrics on the complete question-answering workflow. This end-to-end approach evaluates the entire pipeline with finetuned models for both entity/property extraction and SPARQL generation working together within the multi-agent architecture. Dataset generation quality is evaluated through SPARQL syntax validation, query execution verification, and diversity analysis. Comparative evaluation between the original and enhanced approaches ensures that improvements are properly quantified across different query types and complexity levels.
    \item \textbf{Communication}: The enhanced architecture, implementation details, and experimental findings are documented in this research report. Additionally, the multi-agent system includes comprehensive process visualization, database logging, and semantic logging capabilities that provide enhanced transparency into the decision-making process, allowing users to review and audit system operations through saved execution traces and interactive diagrams. All code, datasets, and models are made available to ensure reproducibility and enable further research in open-source GraphRAG systems.
\end{enumerate}

\section{Experiment Details}
\label{exp-detail}
In general, we will be experimenting with an enhanced GraphRAG system that builds upon the original pipeline through three main improvements. The enhanced system integrates specialized model finetuning for both entity property extraction and SPARQL generation, comprehensive dataset generation, and a sophisticated multi-agent architecture. Additionally, the scope of our enhanced system extends beyond the original framework to provide more robust question-answering capabilities with intelligent strategy selection and adaptive processing paths, while significantly enhancing transparency through comprehensive process visualization, database logging, and reviewable system diagrams. The details are as follows.

\subsection{Overall Architecture Overview}
The enhanced GraphRAG system operates through a coordinated workflow that integrates our three main improvements into a cohesive question-answering pipeline. The system begins with a multi-agent orchestration layer that intelligently routes incoming questions through specialized processing agents. These agents leverage finetuned models for improved entity property extraction and SPARQL generation, and utilize enhanced datasets for training and evaluation purposes.

The overall workflow follows this pattern: (1) a user question enters the multi-agent system through the translation and entity extraction agents, (2) the strategy selection agent determines the optimal processing path based on query complexity and user settings, (3) the system either follows a direct verbalization approach for simple queries or generates SPARQL queries using finetuned models for complex questions, with verbalization failure automatically falling back to SPARQL generation, (4) entity and property search agents utilize the comprehensive generated datasets to improve entity and property matching, and (5) the answer generation agent synthesizes the final response with appropriate fallback mechanisms.

This integrated approach addresses the limitations of the previous adaptive processing system by providing enhanced routing mechanisms, specialized model capabilities for both entity property extraction and SPARQL generation, and comprehensive knowledge coverage through enhanced dataset generation.

\subsection{Original Pipeline (Baseline)}
To establish a proper comparison baseline, we maintain the original GraphRAG framework as described in the previous work \citep{ongris2024frog}. The original system consists of two pipeline versions: a prototype pipeline (v1) focused specifically on Wikidata with limited property scope, and a refined pipeline (v2) that supports multiple knowledge bases but lacks sophisticated reasoning capabilities.

The original pipeline v1 (Figure \ref{fig:pipeline_v1}) serves as our initial baseline, with steps including entity extraction and linking, SPARQL query generation, and answer generation using a predefined list of top-100 frequently used Wikidata properties. The original pipeline v2 (Figure \ref{fig:pipeline_v2}) provides our main baseline comparison, incorporating entities extraction and linking, retrieval mechanisms, and answer generation with support for Wikidata, DBpedia, and local KGs.

\begin{figure}[htbp]
    \centering
    \includegraphics[width=1\textwidth]{assets/pdfs/pipeline-ori.pdf}
    \caption{The Initial Pipeline (Pipeline v1) \citep{ongris2024frog}}
    \label{fig:pipeline_v1}
\end{figure}

\begin{figure}[htbp]
    \centering
    \includegraphics[width=1\textwidth]{assets/pdfs/pipeline2-cropped.pdf}
    \caption{The Upgraded Pipeline (Pipeline v2) \citep{ongris2024frog}}
    \label{fig:pipeline_v2}
\end{figure}

However, both original pipelines are limited by their fixed workflow architecture lacking runtime configuration options, absence of expanded fallback mechanisms, limited operational transparency, lack of specialized model training, and constrained dataset generation capabilities.

\subsection{Enhanced Pipeline Architecture}
Departing from the limitations of the original pipelines, we develop a significantly enhanced architecture through our three main improvements. The enhanced system, implemented as the FrOG (Framework of Open GraphRAG) system, incorporates a sophisticated multi-agent workflow that can adaptively process questions of varying complexity.

The enhanced pipeline operates through the following conceptual components:

\begin{enumerate}
    \item \textbf{Multi-Agent Orchestration}: The system employs a LangGraph-based workflow that coordinates multiple specialized agents, each responsible for specific aspects of the question-answering process. This includes translation agents for multilingual support, entity extraction agents using finetuned models, strategy selection agents for intelligent routing, and specialized processing agents for different query types.
    
    \item \textbf{Adaptive Strategy Selection}: Unlike the fixed processing paths of the original system, our enhanced architecture intelligently determines the optimal approach based on question complexity. Simple questions follow a direct verbalization path, while complex multi-hop queries utilize the sophisticated SPARQL generation pipeline with finetuned models. In cases where verbalization fails for simple questions, the system automatically falls back to SPARQL generation to ensure robust query processing.
    
    \item \textbf{Enhanced Model Integration}: The system incorporates finetuned Large Language Models specifically trained on both entity property extraction and NL2SPARQL tasks using our comprehensive generated datasets. These models provide improved accuracy in both entity and property identification as well as SPARQL query generation, offering better reasoning capabilities for complex knowledge graph queries.
    
    \item \textbf{Comprehensive Dataset Utilization}: The enhanced pipeline leverages both template-based and pattern-based generated datasets for training, evaluation, and property enhancement. This provides broader coverage of query types and improved handling of diverse knowledge domains beyond the limitations of predefined property lists.
    
    \item \textbf{Robust Fallback Mechanisms}: The system includes multiple levels of fallback processing, from verbalization failure to SPARQL generation, and from knowledge graph querying to web search when enabled, ensuring robust performance across different query scenarios.
    
    \item \textbf{Comprehensive Logging and Traceability}: The system implements detailed semantic logging of all processing steps, capturing agent operations, decisions, and reasoning paths in RDF format. This provides complete traceability of the question-answering process, enabling verification, debugging, and transparency into how answers are derived.
\end{enumerate}

The enhanced architecture supports multiple knowledge bases including Wikidata, Curriculum Knowledge Graph, Legal Knowledge Graph, and GESIS Knowledge Graph, while providing significantly improved processing capabilities, broader query coverage, and more reliable performance through its integrated multi-component design.

For model selection, we utilize multiple finetuned versions of state-of-the-art models. The LLMs we use for the enhanced pipeline are \texttt{Mistral NeMo}, \texttt{LLaMA 3.1}, \texttt{Qwen2.5 Coder}, and \texttt{Qwen2.5}. Likewise, all models are the instruction fine-tuned versions. The multi-agent system supports configurable model providers for different processing tasks. The evaluation encompasses both the original Jaccard similarity metrics and additional comprehensive metrics for component-wise analysis.

Finally, we conduct our evaluation using Kaggle's cloud-based infrastructure to ensure computational scalability and reproducibility of our enhanced approach compared to the original baseline systems.

% =============================================================================

\section{Knowledge Bases} \label{knowledgebase}
Our enhanced GraphRAG system is designed to work across multiple knowledge domains to demonstrate its versatility and adaptability. We utilize four distinct knowledge bases that represent different types of structured information:

\begin{itemize}
    \item A large-scale public knowledge graph (Wikidata)
    \item A domain-specific curriculum knowledge graph
    \item A legal document knowledge graph
    \item A scholarly research knowledge graph
\end{itemize}

Each knowledge base presents unique challenges and characteristics that allow us to comprehensively evaluate our enhanced approach across diverse domains and query complexities.

Originally, we planned to include DBpedia as one of our evaluation knowledge bases to demonstrate cross-platform compatibility. However, during our implementation phase, we discovered that the DBpedia Lookup service\footnote{\url{https://lookup.dbpedia.org/}} was unavailable, preventing us from implementing the necessary entity retrieval functionality. Figure~\ref{fig:dbpedia-down} shows the error encountered when attempting to access the service. It should be noted that while the service was down during the development process of this thesis, it became operational again on June 6, 2025. Consequently, we focused our evaluation on the four knowledge bases described below: Wikidata, Curriculum Knowledge Graph, Legal Knowledge Graph, and GESIS Knowledge Graph.

\begin{figure}[htbp]
    \centering
    \includegraphics[width=1\linewidth]{assets/pics/dbpedia-down.png}
    \caption{Screenshot showing DBpedia Lookup service unavailability}
    \label{fig:dbpedia-down}
\end{figure}

\subsection{Wikidata}

Wikidata, originally launched in October 2012 by Wikimedia Foundation \citep{vrandevcic2014}, is an openly accessible, collaborative knowledge base that is editable by anyone \citep{waagmeester2020}. Consequently, the scope of knowledge covered by Wikidata is nearly boundless, just like its sister project, Wikipedia\footnote{\url{https://www.wikipedia.org/}}. However, in contrast to Wikipedia which provides knowledge primarily in the free text format, Wikidata represents information in the structured format \citep{waagmeester2020}. 

The core structure of Wikidata is made up of items, properties, and values, which are connected in statements that closely resemble RDF triples \citep{Cantallops2019ASL}. The knowledge in Wikidata can be retrieved by SPARQL querying through Wikidata Query Service (WDQS)\footnote{\url{https://query.wikidata.org/}}, whose interface can be seen in Figure~\ref{fig:wdqs-ui}. The query example used in Figure~\ref{fig:wdqs-ui} returns the URI of the instances (\texttt{wdt:P31}) of Cat (\texttt{wd:Q146}) alongside with its label.

\begin{figure}[htbp]
    \centering
    \includegraphics[width=1\linewidth]{assets//pics/wdqs.png}
    \caption{Wikidata Query Service's Interface}
    \label{fig:wdqs-ui}
\end{figure}

For our enhanced GraphRAG system, Wikidata serves as the primary large-scale knowledge graph for demonstrating scalability and handling complex multi-hop reasoning scenarios. The vast scale of Wikidata, with millions of entities and thousands of properties, provides an excellent testbed for our enhanced dataset generation methods and finetuned models. Our system utilizes a curated collection of 11,830 Wikidata properties stored in our Weaviate vector database to enable efficient property retrieval and matching during question processing.

For entity retrieval and matching in our enhanced GraphRAG system, we utilize the Wikidata API directly rather than storing all Wikidata entities in our Weaviate vector database. This design decision is driven by the massive scale of Wikidata, which contains 117,632,118 entities as of May 30, 2025\footnote{\url{https://wikidata.org/wiki/Wikidata:Main_Page}}. Storing and indexing this vast number of entities in a vector database would be computationally prohibitive and unnecessary for our use case. Instead, we leverage the Wikidata API's efficient search capabilities for real-time entity retrieval, while maintaining only the 11,830 Wikidata properties in our Weaviate vector database for efficient property matching during question processing.


\subsection{Curriculum Knowledge Graph}
\label{curriculum-kg}

The Curriculum Knowledge Graph represents a domain-specific knowledge base built from Fasilkom UI's Computer Science Major curriculum. This knowledge graph is extracted from the Fasilkom UI's 2024 Curriculum Guidebook for Computer Science Major, which is publicly available through Fasilkom UI's LMS\footnote{\url{https://scele.cs.ui.ac.id/pluginfile.php/1279/block_html/content/buku_panduan_kur_2024_ilmu_komputer.pdf}}. Specifically, the content is obtained from pages 122 to 186, covering the syllabus of all courses.

The curriculum knowledge graph is constructed by extracting entities from the unstructured text data using GLiNER \citep{zaratiana2023glinergeneralistmodelnamed} and converting them into structured RDF triples. The knowledge graph contains comprehensive information about university courses, including course codes, categories, prerequisites, credit values, evaluation methods, and research group associations.

\subsubsection{Knowledge Graph Schema}
The curriculum knowledge graph utilizes a custom ontology with the prefix \texttt{ns1:} representing \texttt{http://example.org/}. The schema includes 8 main properties that capture various aspects of course information, as shown in Table~\ref{table-curriculum-properties}.

\begin{table}[htbp]
\centering
\renewcommand{\arraystretch}{1.15}
\setlength{\tabcolsep}{12pt}
\begin{adjustbox}{width=\textwidth}
    \begin{tabular}{| m{0.4\textwidth} | m{0.6\textwidth} |}
    \hline
    \textbf{Property} & \textbf{Description} \\
    \hline
    \texttt{ns1:has\_course\_code} & Links courses to their unique identifiers \\
    \hline
    \texttt{ns1:has\_course\_category} & Links courses to their classification categories \\
    \hline
    \texttt{ns1:also\_known\_as} & Alternative names or abbreviations for courses \\
    \hline
    \texttt{ns1:has\_prerequisite\_course} & Links courses to their required prerequisites \\
    \hline
    \texttt{ns1:has\_credits} & Links courses to their credit values \\
    \hline
    \texttt{ns1:has\_evaluation\_method} & Links courses to their assessment techniques \\
    \hline
    \texttt{ns1:has\_research\_group} & Links courses to associated research groups \\
    \hline
    \texttt{ns1:has\_minimum\_credits} & Links courses to minimum credit requirements \\
    \hline
    \end{tabular}
\end{adjustbox}
\caption{Curriculum Knowledge Graph Properties}
\label{table-curriculum-properties}
\end{table}

The knowledge graph also defines four main classes: \texttt{ns1:course}, \texttt{ns1:evaluation}, \texttt{ns1:course\_category}, and \texttt{ns1:research\_lab}, which provide the structural foundation for organizing curriculum information. For efficient processing in our enhanced GraphRAG system, we maintain 108 entities and 8 properties in our Weaviate vector database for semantic search and retrieval.

\subsubsection{Sample Data}
The following code shows sample triples about a ``Basis Data" (Database) course in RDF Turtle format:

\lstinputlisting[caption=Sample RDF Triples for Curriculum Knowledge Graph, label=code:curriculum-sample, style=uistyle, language=turtle]{assets/codes/sample-triple.ttl}

The curriculum knowledge graph contains 874 RDF triples in total, providing rich information about course relationships, prerequisites, and academic structure that enables complex reasoning about curriculum pathways and course dependencies.

\subsection{Legal Knowledge Graph}
\label{legal-kg}

The Legal Knowledge Graph is based on the Lex2KG framework \citep{abdurahman2021lex2kg}, which converts Indonesian legal documents from PDF format into structured RDF triples. This knowledge graph contains 784 Indonesian laws (Undang-Undang) with a total size of over 1.1 million triples, making it one of the largest structured legal knowledge bases for Indonesian legislation.

\subsubsection{Knowledge Graph Structure and Modification}
The legal knowledge graph captures various types of legal information including metadata, document structures, textual content, and relationships between legal resources. Table~\ref{table-legal-statistics} shows the distribution of information types in the knowledge graph before our modifications.

\begin{table}[htbp]
\centering
\renewcommand{\arraystretch}{1.15}
\setlength{\tabcolsep}{12pt}
\begin{tabular}{| l | r |}
\hline
\textbf{Information Type} & \textbf{Number of Triples} \\
\hline
Metadata & 6,869 \\
Structure & 714,888 \\
Text & 152,301 \\
Citations in the same law & 22,829 \\
Citations between laws & 1,401 \\
Citations to non-existing resources & 7,095 \\
Amendments & 1,252 \\
Class assignments & 256,416 \\
\hline
\textbf{Total} & \textbf{1,163,051} \\
\hline
\end{tabular}
\caption{Legal Knowledge Graph Statistics (Before Modification)}
\label{table-legal-statistics}
\end{table}

To improve data quality and consistency for our enhanced GraphRAG system, we applied several modifications to the original Lex2KG dataset:

\begin{itemize}
    \item \textbf{Location Standardization}: Fixed inconsistent city names in the \texttt{lex2kg-o:disahkanDi} property (e.g., ``Jakart", ``Jakar" → ``Jakarta")
    \item \textbf{URI Cleanup}: Removed unnecessary ``/text" suffixes from URIs and moved textual content properties to their appropriate parent entities
    \item \textbf{Person Name Standardization}: Normalized variations in person names in the \texttt{lex2kg-o:disahkanOleh} property (e.g., ``JOKO WIDODO" → ``Joko Widodo")
    \item \textbf{Position Title Correction}: Standardized official position titles in the \texttt{lex2kg-o:jabatanPengesah} property
    \item \textbf{Structural Simplification}: Connected entities directly without intermediate ``daftar" (list) entities to reduce query complexity
    \item \textbf{Orphaned Entity Removal}: Removed entities that only appeared as subjects in type declarations without meaningful relationships
\end{itemize}

These modifications resulted in a cleaner, more consistent knowledge graph that better supports our enhanced dataset generation and query processing capabilities. Our vector database contains 243,708 legal entities and 28 legal properties for efficient semantic search.

\subsubsection{Examples of Inconsistencies Before Modification}
We found several inconsistencies in the original legal knowledge graph that resulted from extraction errors during the knowledge graph creation process and required standardization. Figures~\ref{fig:legal-position-query1} and~\ref{fig:legal-position-query2} show SPARQL query examples to retrieve position titles (\texttt{jabatanPengesah}) from the legal knowledge graph, where position titles were inconsistently extracted and formatted as ``PRESIDEN REPUBLIK INDONESIA" (without comma) and ``PRESIDEN REPUBLIK INDONESIA," (with comma) respectively due to variations in the extraction process.
\begin{figure}[htbp]
    \centering
    \includegraphics[width=1\linewidth]{assets/pics/legal-position-query1.png}
    \caption{SPARQL query showing position title formatted as ``PRESIDEN REPUBLIK INDONESIA"}
    \label{fig:legal-position-query1}
\end{figure}
\begin{figure}[htbp]
    \centering
    \includegraphics[width=1\linewidth]{assets/pics/legal-position-query2.png}
    \caption{SPARQL query showing position title formatted as ``PRESIDEN REPUBLIK INDONESIA," with comma}
    \label{fig:legal-position-query2}
\end{figure}
Figure~\ref{fig:legal-document-sample} shows a sample of an Indonesian law document signed by President Joko Widodo, where the document uses ``PRESIDEN REPUBLIK INDONESIA" in all capital letters. The extraction process captured these formatting variations inconsistently, leading to multiple representations of the same entity. These extraction-induced inconsistencies were standardized to ``Presiden Republik Indonesia" in our modified knowledge graph.
\begin{figure}[htbp]
    \centering
    \includegraphics[width=0.8\linewidth]{assets/pics/legal-document-sample.png}
    \caption{Sample of Indonesian law document showing all-caps formatting of position titles}
    \label{fig:legal-document-sample}
\end{figure}
These modifications resolved the extraction errors and resulted in a cleaner, more consistent knowledge graph that better supports our enhanced dataset generation and query processing capabilities. Our vector database contains 243,708 legal entities and 28 legal properties for efficient semantic search.

\subsubsection{Legal Domain Properties}
Key properties in the modified legal knowledge graph include:

\begin{itemize}
    \item \texttt{lex2kg-o:tentang} - Links legal documents to their titles or subject matter
    \item \texttt{lex2kg-o:disahkanPada} - Links documents to their enactment dates
    \item \texttt{lex2kg-o:disahkanOleh} - Links documents to their enacting authorities
    \item \texttt{lex2kg-o:pasal} - Links laws to their articles
    \item \texttt{lex2kg-o:bab} - Links laws to their chapters
    \item \texttt{lex2kg-o:mengubah} - Links amendments to the documents they modify
    \item \texttt{lex2kg-o:merujuk} - Links document sections to what they reference
\end{itemize}

This structure enables complex legal reasoning, such as finding all laws that amend a specific regulation, tracing citation networks between legal documents, and analyzing the evolution of legislation over time.

\subsection{GESIS Knowledge Graph}
\label{gesis-kg}

The GESIS Knowledge Graph (GESIS KG) represents metadata of scientific resources available in the GESIS Search and their semantic relationships in an integrated and consistent form \citep{biswas2025gesis}. GESIS is the Leibniz Institute for the Social Sciences, and their knowledge graph contains comprehensive information about social science research data, publications, and related scholarly resources.

\subsubsection{Knowledge Graph Content and Processing}
The GESIS Knowledge Graph comprises several types of scientific resources and entities, as shown in Table~\ref{table-gesis-resources}.

\begin{table}[htbp]
\centering
\renewcommand{\arraystretch}{1.15}
\setlength{\tabcolsep}{12pt}
\begin{tabular}{| l | l | r |}
\hline
\textbf{Resource Type} & \textbf{Class} & \textbf{Count} \\
\hline
\multicolumn{3}{|c|}{\textbf{Scientific Resources}} \\
\hline
Publications & schema:ScholarlyArticle & 583,085 \\
Datasets & schema:Dataset & 7,546 \\
Instruments & ddi:Instrument & 532 \\
Variables & ddi:Variable & 1,395,499 \\
\hline
\multicolumn{3}{|c|}{\textbf{Entities in Metadata}} \\
\hline
Persons & schema:Person & 466,265 \\
Organizations & schema:Organization & 20,533 \\
Locations & schema:Place & 28,314 \\
Keywords/Concepts & schema:DefinedTerm & 18,297 \\
\hline
\textbf{Total Scientific Resources} & & \textbf{1,986,662} \\
\hline
\end{tabular}
\caption{GESIS Knowledge Graph Resources (Before Modification)}
\label{table-gesis-resources}
\end{table}

\subsubsection{Dataset Modifications}
For our enhanced GraphRAG system, we applied quality improvements to the GESIS dataset to ensure better performance:

\begin{itemize}
    \item \textbf{Duplicate Removal}: Eliminated entities of type \texttt{schema:DuplicateMetadata} to avoid redundant information
    \item \textbf{Incomplete Entity Filtering}: Removed \texttt{schema:ScholarlyArticle} entities that lack essential \texttt{schema:name} properties, ensuring all publications in our dataset have proper identifiers
\end{itemize}

After processing, our enhanced dataset contains 358,530 entities and 33 properties stored in our Weaviate vector database for efficient semantic search and property matching.

\subsubsection{Example of a GESIS Entity}
Figure~\ref{fig:gesis-entity-example} shows an example of a GESIS entity (\texttt{gesis-bib-43272}) as accessible through the GESIS knowledge graph interface. This entity represents a scholarly article titled ``The topicality of the problem of the person in Islamic thought" published in the International Social Science Journal.

\begin{figure}[htbp]
    \centering
    \includegraphics[width=1\linewidth]{assets/pics/gesis-entity-example.png}
    \caption{Example of a GESIS entity (\texttt{gesis-bib-43272}) with its properties and values}
    \label{fig:gesis-entity-example}
\end{figure}

The entity contains various properties including author information (\texttt{Arkoun\_Mohamed}), publication details, language (English), and library location. We can examine this entity using SPARQL queries through the GESIS knowledge graph interface.

A detailed examination of this entity can be performed through different SPARQL queries. Figure~\ref{fig:gesis-entity-subject} shows all properties where \texttt{gesis-bib-43272} is the subject, revealing its attributes such as title, language, and publication information.

\begin{figure}[htbp]
    \centering
    \includegraphics[width=1\linewidth]{assets/pics/gesis-entity-subject.png}
    \caption{SPARQL query results showing properties where \texttt{gesis-bib-43272} is the subject}
    \label{fig:gesis-entity-subject}
\end{figure}

Additionally, Figure~\ref{fig:gesis-entity-object} shows relationships where \texttt{gesis-bib-43272} is the object, demonstrating how this scholarly article is referenced by other entities in the knowledge graph.

\begin{figure}[htbp]
    \centering
    \includegraphics[width=1\linewidth]{assets/pics/gesis-entity-object.png}
    \caption{SPARQL query results showing triples where \texttt{gesis-bib-43272} is the object}
    \label{fig:gesis-entity-object}
\end{figure}

\subsubsection{Scholarly Domain Schema}
The GESIS Knowledge Graph uses established standards and vocabularies, such as schema.org and DDI (Data Documentation Initiative), to ensure interoperability. Key properties include:

\begin{itemize}
    \item \texttt{schema:author} - Links publications to their authors
    \item \texttt{schema:publisher} - Links publications to their publishers
    \item \texttt{schema:datePublished} - Publication dates
    \item \texttt{schema:about} - Links publications to their topics
    \item \texttt{schema:citation} - Citation relationships between publications
    \item \texttt{gesiskg:libraryLocation} - Physical location of resources
    \item \texttt{schema:keywords} - Keywords associated with publications
\end{itemize}

\subsubsection{Link Statistics}
The GESIS Knowledge Graph contains extensive linking between resources, enabling complex scholarly reasoning. Table~\ref{table-gesis-links} shows the distribution of different types of links before our modification process.

\begin{table}[htbp]
\centering
\renewcommand{\arraystretch}{1.15}
\setlength{\tabcolsep}{12pt}
\begin{tabular}{| l | r |}
\hline
\textbf{Link Type} & \textbf{Count} \\
\hline
Automatically generated links between publications and datasets & 78,733 \\
Manually curated links between publications and datasets & 49,509 \\
Links between publications and instruments & 5,813 \\
Links between datasets and variables & 1,393,842 \\
Links between publications & 313,899 \\
Links between publications and variables & 2,954 \\
\hline
\textbf{Total Links} & \textbf{1,861,964} \\
\hline
\end{tabular}
\caption{GESIS Knowledge Graph Link Statistics (Before Modification)}
\label{table-gesis-links}
\end{table}

The GESIS Knowledge Graph serves as an excellent testbed for scholarly domain reasoning, enabling queries about research collaboration networks, publication trends, dataset usage patterns, and the relationships between different types of scholarly resources.

\subsection{Multi-Domain Knowledge Graph Integration}

Our enhanced GraphRAG system's ability to work across these four diverse knowledge domains---general world knowledge (Wikidata), academic curriculum (Curriculum KG), legal documents (Legal KG), and scholarly research (GESIS KG)---demonstrates the versatility and domain-agnostic nature of our approach. Each knowledge graph presents unique challenges:

\begin{itemize}
    \item \textbf{Scale Variation}: From 874 triples (Curriculum) to over 1 million triples (Legal, GESIS)
    \item \textbf{Schema Diversity}: Different ontologies and property patterns across domains
    \item \textbf{Query Complexity}: Varying levels of reasoning complexity from simple property lookups to complex multi-hop reasoning with aggregations
    \item \textbf{Domain Specificity}: From general-purpose world knowledge to highly specialized legal and academic domains
    \item \textbf{Vector Database Integration}: All knowledge graphs are supported by Weaviate collections for efficient semantic search, with collection sizes ranging from 108 entities (University) to 358,530 entities (GESIS)
\end{itemize}

This diversity enables comprehensive evaluation of our enhanced dataset generation methods, finetuned models, and multi-agent architecture across different knowledge representation challenges and domain-specific reasoning requirements. The careful data cleaning and quality improvements applied to each knowledge graph ensure that our enhanced GraphRAG system operates on high-quality, consistent data foundations.

\section{Dataset Generation Methodology}
To overcome the limitations of existing GraphRAG systems that rely on limited dataset coverage and diversity, we develop a comprehensive methodology for creating high-quality Natural Language to SPARQL (NL2SPARQL) training datasets. Our approach differs based on the knowledge graph:

For domain-specific knowledge graphs (Curriculum KG, Legal KG, and GESIS KG), we employ two complementary strategies: template-based generation and random-walk generation to create datasets from scratch. These methods enable us to generate diverse question-SPARQL pairs with varying complexity levels.

For Wikidata, we enhance the existing QALD-9-Plus dataset by augmenting it with: (1) step-by-step reasoning chains (thoughts) generated using Large Language Models, (2) entity matches retrieved from the Wikidata API, and (3) property matches from our Weaviate vector database. This enhancement process adds crucial training signals for entity/property extraction and reasoning, while preserving the high-quality question-SPARQL pairs from QALD-9-Plus.

\subsection{Dataset Enhancement Strategy}
Our dataset preparation follows two distinct approaches based on data availability:

\subsubsection{QALD-9-Plus Enhancement for Wikidata}
For Wikidata, we leverage the established QALD-9-Plus dataset containing 558 validated question-SPARQL pairs. We enhance this dataset through the following process:

\begin{enumerate}
    \item \textbf{Thoughts Generation}: For each question-SPARQL pair, we use Large Language Models to generate step-by-step reasoning chains that explain how to translate the natural language question into the corresponding SPARQL query.
    
    \item \textbf{Entity Enhancement}: We extract entities from both questions and SPARQL queries, then use the Wikidata API to retrieve matching entity URIs and labels. The decision to use the Wikidata API directly, rather than pre-storing entities in our vector database, is driven by Wikidata's massive scale of over 117 million entities. We intentionally add noise entities to improve model robustness during training.
    
    \item \textbf{Property Enhancement}: We identify properties used in SPARQL queries and retrieve matching properties from our Weaviate vector database containing 11,830 Wikidata properties. Similar to entities, we add noise properties to enhance training diversity.
    
    \item \textbf{Structured Output Format}: Each enhanced entry includes the original question and SPARQL, generated thoughts, entity matches (including noise), and property matches (including noise).
\end{enumerate}

This enhancement process transforms QALD-9-Plus into a richer training dataset suitable for fine-tuning both entity/property extraction and SPARQL generation models.

\subsubsection{Full Generation for Domain-Specific Knowledge Graphs}
For Curriculum KG, Legal KG, and GESIS KG, where no existing evaluation datasets exist, we generate complete datasets from scratch using template-based and random-walk approaches with complexity stratification, as detailed in the following sections.

\begin{figure}
    \centering
    \includegraphics[width=0.7\linewidth]{assets/pics/dataset1.png}
    \caption{Overview of the dual dataset generation methodology combining template-based and random-walk approaches}
    \label{fig:enter-label}
\end{figure}

\subsection{Template-Based Dataset Generation}
The template-based dataset generation methodology follows a discovery-first approach where predefined question templates are instantiated with valid entity-property combinations discovered through graph exploration. This ensures that generated questions have corresponding valid and executable SPARQL queries that return meaningful results from the knowledge graph.

\subsubsection{Template Structure}
Each template contains the following components:

\begin{enumerate}
    \item \textbf{Question Templates}: Multiple natural language formulations of a question type, allowing for linguistic diversity. For multilingual support, templates can include variations in different languages.
    
    \item \textbf{SPARQL Template}: A parameterized SPARQL query pattern with placeholders for entities and properties that will be discovered from the knowledge graph.
    
    \item \textbf{Chain of Thoughts}: Step-by-step reasoning explanations that describe the process of translating the natural language question to a SPARQL query, providing valuable training signals for models.
    
    \item \textbf{Complexity Classification}: Each template is categorized by complexity level (basic, intermediate, advanced) to ensure balanced coverage across difficulty levels.
    
    \item \textbf{Category}: Domain classification (e.g., scholarly, legal) to organize templates by knowledge domain.
\end{enumerate}

Templates are designed with varying complexity levels to provide comprehensive coverage:

\begin{itemize}
    \item \textbf{Basic Templates (40\%)}: Simple one-hop relationships focusing on direct subject-predicate-object patterns.
    \item \textbf{Intermediate Templates (30\%)}: Two-hop relationships and queries with aggregation functions like COUNT.
    \item \textbf{Advanced Templates (30\%)}: Complex multi-hop relationships with filters, nested patterns, and sophisticated aggregations.
\end{itemize}

\begin{lstlisting}[
  language=json, 
  caption=Example of a Generated Dataset Entry with Formatted SPARQL, 
  label=lst:dataset_example,
  breaklines=true,
  breakatwhitespace=true,
  breakanywhere=true,
  postbreak=\mbox{\textcolor{red}{$\hookrightarrow$}\space},
  basicstyle=\footnotesize\ttfamily,
  columns=flexible,
  keepspaces=true,
  showstringspaces=false,
  frame=single,
  xleftmargin=0pt,
  xrightmargin=0pt,
  linewidth=\linewidth,
  resetmargins=true,
  breaklines=true
]
{
  "id": "q1",
  "question": "How many credits does the Algorithm Design and Analysis course have?",
  "sparql": "SELECT?value\nWHERE {\n  ns1:mobile_technology ns1:has_credits?value.\n}",
  "category": "university",
  "complexity": "basic",
  "templateId": "course-credits",
  "thoughts": [
    "1. The question asks for the credit value of Algorithm Design and Analysis.",
    "2. The entity 'Algorithm Design and Analysis' represents a course in the university domain.",
    "3. The property 'has_credits' links a course to its credit value.",
    "4. To solve this, retrieve the credit value linked to Algorithm Design and Analysis via the 'has_credits' property.",
    "5. Construct a SPARQL query to retrieve the credit value for Algorithm Design and Analysis."
  ],
  "entities": ["Algorithm Design and Analysis"],
  "properties": ["has_credits"],
  "entities_matches": [
    {"id": "http://example.org/algorithm_design_and_analysis", 
     "label": "Algorithm Design and Analysis"}
  ],
  "properties_matches": [
    {"id": "http://example.org/has_credits", 
     "label": "has credits"}
  ]
}
\end{lstlisting}

% Place the template-based approach workflow diagram here
% ```mermaid
% flowchart TD
%     A[Template Definition] --> B[Placeholder Extraction]
%     B --> C[Discovery Query Generation]
%     C --> D[KG Query Execution]
%     D --> E[Valid Entity-Property Selection]
%     E --> F[Template Instantiation]
%     F --> G[Question + SPARQL Generation]
%     G --> H[Chain of Thoughts Generation]
%     H --> I[Entity-Property Matching]
% ```

\begin{figure}
    \centering
    \includegraphics[width=0.75\linewidth]{assets/pics/template2.jpg}
    \caption{Workflow of the template-based dataset generation process showing the discovery-first approach and SPARQL post-processing}
    \label{fig:enter-label}
\end{figure}

\subsubsection{Discovery-First Instantiation Process}
The key innovation in our template-based approach is the discovery-first methodology that guarantees valid instantiations:

\begin{enumerate}
    \item \textbf{Placeholder Extraction}: Identify all placeholders in the template (e.g., \texttt{\{entity\}}, \texttt{\{property\}}).
    
    \item \textbf{Discovery Query Generation}: Transform the SPARQL template into a discovery query that finds all valid combinations of entities and properties that would yield results in the knowledge graph.
    
    \item \textbf{Validation and Selection}: Execute the discovery query against the knowledge graph and randomly select valid combinations from the results.
    
    \item \textbf{Template Instantiation}: Replace placeholders in both the question and SPARQL templates with the selected entities and properties.
    
    \item \textbf{Chain of Thoughts Generation}: Generate step-by-step reasoning using the thoughts template, with context-aware replacement of URIs and labels.
\end{enumerate}

This discovery-first methodology ensures that all generated question-SPARQL pairs are valid and executable, eliminating the need for post-generation validation and filtering.

\subsubsection{Entity and Property Matching}
To enhance the dataset's utility for training entity linking and property matching components, our methodology incorporates a sophisticated matching system:

\begin{enumerate}
    \item \textbf{Vector Embedding}: All entities and properties in the knowledge graph are encoded using sentence transformer models to create semantic vector representations.
    
    \item \textbf{N-gram Analysis}: Questions are processed into n-grams to capture multi-word expressions and handle cases where properties are not explicitly stated.
    
    \item \textbf{Hybrid Search}: A combination of keyword-based and semantic similarity search identifies relevant entities and properties mentioned in the question.
    
    \item \textbf{Match Recording}: For each question-SPARQL pair, the system records which entities and properties from the SPARQL query were successfully matched in the question text, creating valuable training data for entity and property extraction models.
\end{enumerate}

\subsection{Random-Walk Dataset Generation}
While the template-based approach provides high-quality, linguistically diverse questions, it may not capture the full range of structural patterns present in a knowledge graph. To address this limitation, we developed a pattern-based random walk methodology that discovers valid query patterns directly from the graph structure.
\begin{figure}
    \centering
    \includegraphics[width=0.7\linewidth]{assets/pics/salib2.jpg}
    \caption{Workflow of the random-walk dataset generation process showing pattern discovery, question generation, and SPARQL post-processing}
    \label{fig:enter-label}
\end{figure}

\subsubsection{Pattern Discovery}
The random walk methodology systematically explores the knowledge graph to discover valid query patterns with increasing complexity:
\begin{enumerate}
    \item \textbf{1-Property Pattern (50\%)}: Basic patterns involving a single triple pattern with one fixed entity and one variable:
    \begin{itemize}
        \item Linear Target Pattern:  \texttt{fixedEntity property ?target}
    \end{itemize}
    
    \begin{figure}[H]
        \centering
        \includegraphics[width=0.5\textwidth]{assets/pics/1prop.png}
        \caption{1-Property Pattern Visualization}
        \label{fig:1-property}
    \end{figure}
    
    \item \textbf{2-Property Patterns (30\%)}: Intermediate patterns involving two connected triple patterns:
    \begin{itemize}
        \item Linear Target Pattern: \texttt{fixedEntity prop1 ?hidden . ?hidden prop2 ?target} 
        \item Branching Pattern: \texttt{fixedEntity1 prop1 ?target . ?target prop2 fixedEntity2}
    \end{itemize}
    
    \begin{figure}[H]
        \centering
        \includegraphics[width=0.5\textwidth]{assets/pics/2prop.jpg}
        \caption{2-Property Pattern Visualization}
        \label{fig:2-property}
    \end{figure}
    
    \item \textbf{3-Property Patterns (20\%)}: Advanced patterns involving three connected triple patterns:
    \begin{itemize}
        \item Linear End Pattern: \texttt{entity prop1 ?h1 . ?h1 prop2 ?h2 . ?h2 prop3 ?target}
        \item Linear Middle Pattern: \texttt{entity1 prop1 ?h . ?h prop2 ?target . ?target prop3 entity2}
        \item Star Pattern: \texttt{?hidden prop1 entity1 . ?hidden prop2 entity2 . ?hidden prop3 ?target}
    \end{itemize}
    
    \begin{figure}[H]
        \centering
        \includegraphics[width=0.6\textwidth]{assets/pics/3prop.jpg}
        \caption{3-Property Pattern Visualization}
        \label{fig:3-property}
    \end{figure}
\end{enumerate}

\subsubsection{Validation and Variation Generation}
To ensure diversity and quality, our methodology includes:

\begin{enumerate}
    \item \textbf{Pattern Validation}: Each generated pattern is validated by executing it against the knowledge graph to ensure it returns non-empty results.
    
    \item \textbf{Structural Variations}: For each pattern type, multiple variations are generated by systematically reversing the direction of triple patterns. For example, a 3-property pattern can generate up to $2^3 = 8$ different variations.
    
    \item \textbf{Quality Filtering}: Low-quality patterns involving system properties or very common properties are excluded to ensure meaningful and diverse queries.
\end{enumerate}

\subsubsection{Natural Language Question Generation}
Since the random walk approach generates SPARQL queries without natural language questions, we employ a sophisticated question generation process:

\begin{enumerate}
    \item \textbf{Entity and Property Extraction}: For each generated SPARQL query, we extract the entities and properties used in the query.
    
    \item \textbf{Template Context Transfer}: Using templates from the template-based approach as examples, we create contextual prompts for Large Language Models.
    
    \item \textbf{LLM-Based Generation}: We employ Large Language Models to generate natural language questions and step-by-step reasoning chains for each SPARQL query, ensuring linguistic quality and diversity.
    
    \item \textbf{Quality Control}: Generated questions are filtered to ensure they accurately represent the information need expressed by the SPARQL query.
\end{enumerate}

\subsection{Hybrid Approach Benefits}
The combination of template-based and random-walk methodologies provides several advantages for domain-specific knowledge graphs:

\begin{enumerate}
    \item \textbf{Comprehensive Coverage}: Template-based generation ensures linguistic diversity and question quality, while random-walk discovery captures the full range of structural patterns in the knowledge graph.
    
    \item \textbf{Guaranteed Validity}: Both approaches ensure that generated SPARQL queries are valid and return results when executed against the knowledge graph.
    
    \item \textbf{Balanced Complexity}: The stratified approach to complexity levels ensures a balanced distribution of query types for effective model training.
    
    \item \textbf{Domain Adaptability}: The methodology can be applied to different knowledge domains by adapting templates and discovery parameters to domain-specific requirements.
\end{enumerate}

For Wikidata, while we do not generate new question-SPARQL pairs, our enhancement process adds similar benefits by incorporating thoughts generation, entity/property matching with noise, and structured training formats that align with our generated datasets for other knowledge graphs. This ensures consistency in training data format across all domains while leveraging the high-quality QALD-9-Plus foundation.

\subsection{Dataset Post-Processing and SPARQL Formatting}
\label{sec:dataset-postprocessing}

To enhance the quality of training data and improve model learning efficiency, we apply comprehensive post-processing to all generated datasets across both Wikidata and domain-specific knowledge graphs. Our empirical observations revealed that base language models tend to generate SPARQL queries with specific formatting patterns, including uppercase keywords, proper indentation, and specific spacing conventions. To align our training data with these natural tendencies, we implement a systematic post-processing pipeline.

\subsubsection{SPARQL Query Formatting Pipeline}

The post-processing pipeline consists of four key transformations applied to all SPARQL queries in our datasets:

\begin{enumerate}
    \item \textbf{Keyword Capitalization}: All SPARQL keywords (SELECT, WHERE, FILTER, BIND, UNION, ORDER BY, COUNT, AS, GROUP BY, LIMIT, OFFSET, DISTINCT, ASK) are converted to uppercase to match the conventional SPARQL formatting style.
    
    \item \textbf{URI to Prefix Conversion}: Full URIs are replaced with prefixed notation for improved readability. For example, \texttt{<http://example.org/has\_credits>} becomes \texttt{ns1:has\_credits}, making queries more concise and easier for models to process. However, for Legal and GESIS knowledge graphs, entity URIs are not converted to prefixed notation because many entities in these knowledge graphs contain forward slash characters in their identifiers (e.g., \texttt{lex2kg:uu/1950/1}), which would result in invalid SPARQL syntax when combined with prefix notation. Therefore, these entities remain in their full URI format enclosed in angle brackets.
    
    \item \textbf{Query Structure Formatting}: SPARQL queries are reformatted with proper indentation and line breaks:
    \begin{itemize}
        \item WHERE clauses are placed on separate lines with proper brace alignment
        \item Triple patterns within WHERE clauses are indented consistently
        \item Special clauses (FILTER, BIND, UNION) are formatted on separate lines
        \item ORDER BY, GROUP BY, and LIMIT clauses are properly separated
    \end{itemize}
    
    \item \textbf{Whitespace Normalization}: Spaces before question marks (\texttt{?}) and dots (\texttt{.}) are removed to ensure consistent spacing patterns throughout the queries.
\end{enumerate}

\subsubsection{Example of Post-Processing Transformation}

Before post-processing:
\begin{lstlisting}[language=SPARQL, caption=SPARQL query before post-processing]
select ?value where { <entity> <property> ?value . }
\end{lstlisting}

After post-processing:
\begin{lstlisting}[language=SPARQL, caption=SPARQL query after post-processing]
SELECT?value
WHERE {
  ns1:entity ns1:property?value.
}
\end{lstlisting}

This formatting approach significantly improves model performance by:
\begin{itemize}
    \item Reducing the learning complexity for SPARQL syntax patterns
    \item Providing consistent formatting that aligns with model generation preferences
    \item Enhancing query readability for debugging and validation
    \item Standardizing the representation across different knowledge graph domains
\end{itemize}

The post-processing is applied uniformly to:
\begin{itemize}
    \item Enhanced QALD-9-Plus dataset for Wikidata
    \item Template-based generated datasets for domain-specific KGs
    \item Random-walk generated datasets for domain-specific KGs
\end{itemize}

\begin{figure}[htbp]
    \centering
    \includegraphics[width=1\linewidth]{assets/pics/asdf.jpg}
    \caption{SPARQL Query Post-Processing Pipeline showing the four-stage transformation process}
    \label{fig:sparql-postprocessing}
\end{figure}

% ---------------------------------------

\section{Model Finetuning Methodology}
\label{model-finetune}
To overcome limitations in previous approaches that relied on few-shot prompting, we formulated a specialized model fine-tuning methodology using Quantized Low-Rank Adaptation (QLoRA) for two critical components: entity and property extraction, and SPARQL query generation. This methodological approach significantly improves performance over zero-shot or few-shot approaches by tailoring model weights specifically for knowledge graph tasks while maintaining computational efficiency. We designed distinct hyperparameter configurations for each task based on their complexity and the nature of the required outputs.

\begin{figure}[htbp]
    \centering
        \includegraphics[width=0.8\textwidth]{assets/pics/lora.jpg}
    \caption{QLoRA Architecture: Only small adapter modules are trained while the base model remains frozen, reducing trainable parameters by over 99\% compared to full fine-tuning}
    \label{fig:lora-architecture}
\end{figure}

\subsection{Entity and Property Extraction Model}
\label{entity-prop-extract}
The entity and property extraction model methodology serves as the foundational processing step in our pipeline, designed to identify relevant entities and properties mentioned in natural language questions. This methodology enables accurate mapping between question terms and knowledge graph elements, which is essential for generating correct SPARQL queries.

\begin{figure}[htbp]
    \centering
    \includegraphics[width=0.4\textwidth]{assets/pics/extract.jpg}
    \caption{Entity and Property Extraction Model: Takes natural language questions as input and identifies entities and properties for SPARQL generation}
    \label{fig:entity-extraction}
\end{figure}

\subsubsection{Task Definition and Approach}
Our methodological approach for entity and property extraction focuses on identifying two key elements from natural language questions:
\begin{itemize}
    \item All entities mentioned in the question (e.g., people, organizations, concepts)
    \item All properties relevant to answering the question (e.g., ``author", ``located in", ``published date")
\end{itemize}

We designed the output to follow a structured format with ``entities" and ``properties" keys, each containing an array of relevant terms identified from the question. This structured approach ensures consistent processing in subsequent pipeline stages and allows for clear evaluation of extraction quality.

\subsubsection{Model Selection Criteria and Adaptation Strategy}
Our model selection methodology focused on large language models with demonstrated capabilities in understanding natural language semantics, contextual interpretation, and structured output generation. We evaluated state-of-the-art models including Mistral, Qwen, and Llama families, prioritizing instruction-tuned versions that have been optimized for following specific task instructions.

For the entity and property extraction task, we designed a lighter QLoRA configuration since this task requires less adaptation from the base model's capabilities:

\begin{itemize}
    \item \textbf{Rank (r)}: 16 - Selected as the optimal rank through empirical analysis of extraction quality
    \item \textbf{Alpha ($\alpha$)}: 16 - Matched to the rank parameter for balanced adaptation
    \item \textbf{Target modules}: Selected attention layers and MLP layers for fine-tuning
    \item \textbf{Dropout}: 0 - Chosen to optimize for deterministic behavior in extraction tasks
    \item \textbf{Bias training}: None - Excluded to minimize trainable parameters
    \item \textbf{QLoRA configuration}: Standard approach without rank stabilization, as extraction tasks showed less susceptibility to training instabilities
\end{itemize}

This configuration approach results in approximately 40-60 million trainable parameters (depending on model size), representing less than 1\% of the original model parameters, enabling efficient adaptation while capturing domain-specific extraction patterns.

\begin{figure}[htbp]
    \centering
   \includegraphics[width=0.5\linewidth]{assets/pics/paramcompar.jpg}
    \caption{Comparison of parameter efficiency: LoRA fine-tuning requires less than 1\% of the parameters compared to full model fine-tuning}
    \label{fig:parameter-efficiency}
\end{figure}

\subsubsection{Training Data Methodology}
Our methodology for training data preparation focused on creating diverse, representative examples of entity and property extraction tasks. We designed a comprehensive approach that incorporated data from multiple sources:

\begin{itemize}
    \item For Wikidata: Enhanced QALD-9-Plus dataset with LLM-generated thoughts and entity/property annotations (including noise) retrieved from Wikidata API and Weaviate
    \item For domain-specific KGs: Systematically generated examples from our template-based dataset generation system
    \item For all KGs: Entity-property annotations derived from our datasets generation methodology (domain-specific KGs) or enhancement process (Wikidata)
\end{itemize}

To ensure optimal learning transfer, we structured the training data with a consistent pattern:
\begin{itemize}
    \item System instructions defining the extraction task parameters
    \item Question inputs representing diverse knowledge graph domains
    \item Expected structured outputs with accurately identified entities and properties
\end{itemize}

This structured approach to training data preparation ensured comprehensive coverage of extraction patterns while maintaining consistency in output formatting.

\subsubsection{Training Methodology}
Our training methodology was designed to optimize learning while preventing overfitting. Key methodological decisions included:

\begin{itemize}
    \item \textbf{Batch organization}: Small per-device batches with gradient accumulation to balance memory constraints and effective batch size
    \item \textbf{Learning rate strategy}: Moderate initial rate with linear decay schedule
    \item \textbf{Training duration}: Limited to 5 epochs, determined empirically to prevent overfitting
    \item \textbf{Regularization approach}: Applied weight decay to enhance generalization
    \item \textbf{Numerical precision}: Optimized based on hardware capabilities
    \item \textbf{Validation strategy}: Regular evaluation during training to monitor performance
\end{itemize}

We employed unpacked training to preserve the structure of individual examples and enable better learning of the structured output format. Our model selection criteria prioritized validation loss as the key metric, ensuring optimal performance while preventing overfitting.

\subsection{SPARQL Generation Model}
\label{sparql-gen}
The SPARQL generation model methodology represents the core reasoning component of our system, transforming identified entities and properties into executable SPARQL queries that accurately represent the user's information need. This task requires a more sophisticated methodological approach than entity extraction, as it necessitates understanding of graph query patterns, SPARQL syntax, and domain-specific knowledge graph structures.

\begin{figure}[htbp]
    \centering
    \includegraphics[width=0.4\linewidth]{assets/pics/spaqrgen.jpg}
    \caption{SPARQL Generation Model: Transforms natural language questions and candidate entities/properties into executable SPARQL queries with reasoning steps}
    \label{fig:sparql-generation}
\end{figure}

\subsubsection{Task Definition and Approach}
Our methodological approach for SPARQL generation focused on the transformation of natural language questions and candidate entities/properties into valid, executable SPARQL queries. This process requires:

\begin{itemize}
    \item Selection of appropriate entities and properties from candidate lists
    \item Determination of query structure based on question intent
    \item Generation of syntactically correct SPARQL with domain-appropriate prefixes
    \item Implementation of variable binding, filtering, and result formatting
    \item Incorporation of step-by-step reasoning to explain query construction
\end{itemize}

The training data preparation incorporates our comprehensive SPARQL post-processing pipeline (see Section~\ref{sec:dataset-postprocessing}), ensuring that all training examples follow consistent formatting patterns that align with the natural generation tendencies of language models. This formatting consistency significantly improves the model's ability to generate syntactically correct and properly formatted SPARQL queries.

\subsubsection{Model Adaptation Strategy}
For SPARQL generation, we designed domain-specific adaptation configurations due to the task's complexity and the need for precise syntax generation:

\begin{itemize}
    \item \textbf{Rank (r)}: Varied by knowledge graph domain - 64 for Wikidata and Curriculum KG, 32 for Legal and GESIS KG (which have longer inputs with more complex data structures)
    \item \textbf{Alpha ($\alpha$)}: Matched to the rank parameter for consistent adaptation (64 or 32 depending on domain)
    \item \textbf{Target modules}: Comprehensive selection of attention and MLP layers
    \item \textbf{Dropout}: 0.1 - Increased dropout for enhanced generalization on complex tasks
    \item \textbf{Bias training}: Included to improve syntax pattern learning
    \item \textbf{QLoRA configuration}: Rank-Stabilized approach to prevent training instabilities
\end{itemize}

This methodological approach results in approximately 160-230 million trainable parameters, representing less than 3\% of the original model parameters, but providing sufficient capacity to learn complex SPARQL patterns across different knowledge graph domains.

\subsubsection{Training Data Methodology}
Our SPARQL generation training data methodology focused on creating comprehensive examples that included:

\begin{itemize}
    \item Diverse natural language questions across knowledge domains
    \item Structured candidate entity and property lists with metadata (including intentional noise for robustness)
    \item Step-by-step reasoning chains explaining query construction logic
    \item Syntactically correct target SPARQL queries for each knowledge graph domain
\end{itemize}

All SPARQL queries in our training datasets undergo comprehensive post-processing to align with the natural generation patterns observed in base language models. This includes converting keywords to uppercase, replacing full URIs with prefixed notation, applying consistent indentation and formatting, and normalizing whitespace. This preprocessing step ensures that our fine-tuned models learn from consistently formatted examples that match their inherent generation preferences, leading to more reliable SPARQL query generation. The post-processing pipeline is detailed in Figure~\ref{fig:sparql-postprocessing}.

For Wikidata, we enhanced the QALD-9-Plus dataset by:
\begin{itemize}
    \item Generating thoughts using Large Language Models to explain the reasoning process
    \item Adding entity matches from Wikidata API with noise entities for training robustness
    \item Including property matches from our Weaviate vector database with noise properties
\end{itemize}

For domain-specific knowledge graphs (Curriculum KG, Legal KG, GESIS KG), we used our template-based and random-walk generation approaches to create complete datasets with built-in thoughts and entity/property annotations.

\subsubsection{Training Methodology}
Our SPARQL generation training methodology employed a more intensive approach:

\begin{itemize}
    \item \textbf{Batch organization}: Smaller per-device batches with increased gradient accumulation
    \item \textbf{Learning rate strategy}: Conservative rate with polynomial decay for stability
    \item \textbf{Training duration}: Extended to 10 epochs to accommodate the complexity of the task
    \item \textbf{Regularization approach}: Enhanced weight decay for stronger regularization
    \item \textbf{Sequence handling}: Enabled sequence packing for varied sequence lengths
\end{itemize}

We incorporated advanced methodological optimizations including gradient computation strategies and batch organization techniques to handle the longer sequences involved in SPARQL generation training.

\begin{figure}[htbp]
    \centering
        \includegraphics[width=0.8\linewidth]{assets/pics/training.jpg}
    \caption{Training Methodology: From data preparation to model selection, showing the complete fine-tuning approach}
    \label{fig:training-process}
\end{figure}

\subsection{Cross-Domain Adaptation Methodology}
\label{cross-domain}

A key innovation in our fine-tuning methodology is the cross-domain adaptation approach, which enables efficient development across different knowledge graph domains. We formulated a sequential experimental approach to optimize our methodology:

Our methodology follows these key principles:

\begin{figure}[htbp]
    \centering
    \includegraphics[width=0.8\linewidth]{assets/pics/cross.jpg}
    \caption{Cross-Domain Adaptation Methodology: Initial adaptation on Wikidata followed by domain-specific adaptation to specialized knowledge graphs}
    \label{fig:cross-domain}
\end{figure}

\begin{enumerate}
    \item \textbf{Experimental Foundation}: We first conducted comprehensive experiments using Wikidata as our primary testbed. This phase focused on:
    \begin{itemize}
        \item Testing various prompt structures to determine optimal instruction formats
        \item Exploring different dataset formats and organization methods
        \item Evaluating multiple finetuning parameter configurations (learning rates, LoRA ranks, etc.)
    \end{itemize}
    
    \item \textbf{Methodology Knowledge Transfer}: Our approach transfers the methodological insights gained from Wikidata experiments:
    \begin{itemize}
        \item Each domain-specific knowledge graph (Curriculum KG, Legal KG, GESIS KG) is finetuned directly from the base model
        \item We apply the optimal prompt structures, dataset formats, and hyperparameters discovered in the Wikidata experiments
        \item This ensures each domain model is specialized to its knowledge domain
    \end{itemize}
    
    \item \textbf{Domain-Specific Adaptations}: For each knowledge graph, we implemented specific adaptations:
    \begin{itemize}
        \item Domain-specific prefixes and namespace conventions
        \item Schema-aware prompt structures reflecting the ontology of each domain
        \item Query pattern templates appropriate for each knowledge domain
    \end{itemize}
\end{enumerate}

This methodological approach significantly reduces experimental overhead while ensuring each domain receives optimized finetuning configurations. By using Wikidata as our experimental foundation, we identified best practices that could be directly applied to other domains without requiring extensive parameter searching for each new knowledge graph.

The workflow ensures that each domain-specific model is directly finetuned from the base model, preserving domain purity while benefiting from the methodological insights gained through our Wikidata experiments.

\subsection{Computational Efficiency Methodology}
To enable adaptation of large language models on standard research infrastructure, we developed a comprehensive methodology for computational efficiency:

\begin{figure}[htbp]
    \centering
    \includegraphics[width=0.8\linewidth]{assets/pics/eff.png}
    \caption{Computational Efficiency Methodology: Enabling adaptation of 7-12B parameter models on standard research infrastructure}
    \label{fig:memory-optimization}
\end{figure}

\begin{itemize}
    \item \textbf{Numerical precision methodology}: Adoption of 4-bit precision, reducing memory requirements by up to 8x
    \item \textbf{Parameter-Efficient Adaptation}: Application of QLoRA to reduce trainable parameters to less than 3\% of the full model
    \item \textbf{Gradient computation strategy}: Implementation of gradient checkpointing to optimize memory usage
    \item \textbf{Optimizer methodology}: Utilization of 8-bit optimization to reduce state memory footprint
    \item \textbf{Batch processing approach}: Application of gradient accumulation for effective batch size management
\end{itemize}

This methodological approach to computational efficiency enables adaptation of large models (7-12 billion parameters) on standard research infrastructure, making our approach accessible and reproducible for researchers with limited computational resources.

\begin{figure}[htbp]
    \centering
    \includegraphics[width=0.8\linewidth]{assets/pics/learning2.jpg}
    \caption{Learning Progression: Training and validation loss patterns showing typical convergence behavior during model adaptation}
    \label{fig:learning-curves}
\end{figure}

This comprehensive fine-tuning methodology provides significant advantages over the few-shot prompting approach used in the original GraphRAG system, enabling more accurate entity extraction and SPARQL generation across diverse knowledge graph domains while maintaining computational efficiency through parameter-efficient adaptation techniques.

\subsection{Model Deployment Methodology}
To bridge the gap between model training and deployment in our multi-agent system, we implement a standardized model deployment approach. All fine-tuned models are converted to GGUF\footnote{\url{https://github.com/ggerganov/ggml/blob/master/docs/gguf.md}} (GPT-Generated Unified Format) format, which significantly reduces model size while maintaining inference quality. This conversion process optimizes the models for runtime efficiency with 4-bit quantization, reducing the memory footprint by up to 8x compared to full-precision models.

For model serving within our multi-agent architecture, we utilize Ollama\footnote{\url{https://ollama.com/}}, an open-source runtime environment for efficient local model deployment. Ollama provides a lightweight API interface that allows the LangGraph agents to make inference requests to the deployed models through a consistent REST API, simplifying the integration of multiple fine-tuned models into our production workflow.

This deployment approach creates a clean separation between model training and model serving components, enabling independent scaling and versioning of models within the multi-agent system. Each LangGraph agent accesses the appropriate fine-tuned model (entity extraction or SPARQL generation) through standardized API calls, ensuring consistent model availability across the workflow.

% --------------------------

\section{Multi-Agent Architecture Design}
The third core improvement of our enhanced GraphRAG system is the implementation of a sophisticated multi-agent architecture that enables intelligent routing, adaptive processing, and robust fallback mechanisms. This approach significantly enhances the capabilities of the original framework by replacing the single-agent architecture with a coordinated system of specialized agents that can collaboratively solve complex knowledge graph queries.

\subsection{Agent Workflow Design}
\label{agent-workflow}
\begin{figure}[htbp]
    \centering
    \includegraphics[width=1\linewidth]{assets/pics/multiiii.jpg}
    \caption{Multi-Agent Architecture of the Enhanced GraphRAG System}
    \label{fig:multi-agent-architecture}
\end{figure}
The multi-agent architecture is designed as a directed graph of specialized processing nodes, each responsible for a specific aspect of the question-answering process. The workflow follows the LangGraph paradigm, which enables the construction of complex agent systems with conditional routing and state management.

As illustrated in Figure~\ref{fig:multi-agent-architecture}, the multi-agent system consists of the following specialized agents:

\begin{enumerate}
    \item \textbf{Translation Agent}: The entry point of the workflow responsible for detecting the language of the input question and translating non-English queries to English for processing. This ensures the system can handle multilingual queries while maintaining high accuracy.
    
    \item \textbf{Entity Extraction Agent}: Analyzes the (translated) natural language question to identify relevant entities and properties mentioned. This agent leverages a fine-tuned language model specifically trained for entity and property extraction tasks as described in Section~\ref{entity-prop-extract}. The fine-tuned model is deployed in GGUF format and served through Ollama, allowing efficient integration within the LangGraph workflow.
    
    \item \textbf{Strategy Selection Agent}: Determines the optimal processing approach based on query complexity, extracted entities, and user settings. This critical decision point routes the query to either the verbalization path (for simple queries) or the SPARQL generation path (for complex queries).
    
    \item \textbf{Verbalization Agent}: Handles simple, single-entity queries by directly retrieving property information from the knowledge graph and presenting it in natural language format. This agent is optimized for efficiency with single-hop queries.
    
    \item \textbf{Property Generation Agent}: Enhances the list of potential properties for a query using semantic search techniques against the property vector database, improving the coverage of relevant properties beyond those initially extracted.
    
    \item \textbf{SPARQL Generation Agent}: Generates executable SPARQL queries for complex questions using a fine-tuned language model specifically trained on NL2SPARQL tasks as described in Section~\ref{sparql-gen}. Like the entity extraction model, this fine-tuned model is deployed in GGUF format and accessed through Ollama's API interface, enabling consistent model availability within the multi-agent system.

    \item \textbf{Google Search Agent}: Provides an optional fallback mechanism when knowledge graph approaches fail to provide sufficient results, ensuring the system can still deliver useful information even when faced with queries outside the knowledge graph's domain.
    
    \item \textbf{Answer Generation Agent}: Synthesizes the final answer based on query results, formatting the response in natural language that directly addresses the user's question.
\end{enumerate}

Each agent maintains and updates a shared state object that contains the cumulative processing information, enabling seamless information transfer between components while maintaining a coherent processing context. The architecture supports both synchronous and asynchronous agent operations, with each agent designed to handle specific aspects of the question-answering process.

The workflow operates through a series of transitions between agents, with conditional routing based on the state of the processing pipeline. As shown in Figure~\ref{fig:multi-agent-architecture}, the system begins with the Translation Agent, which passes the processed question to the Entity Extraction Agent. After entities are extracted, the Strategy Selection Agent determines the optimal processing path based on query complexity and user settings.

For simple questions, the Verbalization Agent attempts to provide a direct answer by retrieving property information for the identified entity. For complex questions, or when verbalization fails, the system routes to the Property Generation Agent which enhances the list of potential properties before passing control to the SPARQL Generation Agent. If both verbalization and SPARQL generation fail to produce satisfactory results, the system can optionally fall back to the Google Search Agent if enabled, providing information from external sources outside the knowledge graph. The Answer Generation Agent then synthesizes the final response based on the query results.

The coordination between agents is facilitated through a state object that maintains the cumulative processing context, including:

\begin{itemize}
    \item The original and translated questions
    \item Extracted entities and properties
    \item Intermediate processing results
    \item Generated SPARQL queries
    \item Query execution results
    \item Processing approach used (verbalization, SPARQL, or Google Search)
    \item Runtime settings and configuration
\end{itemize}

This state-based approach ensures that each agent has access to all relevant information from previous processing steps, enabling more informed decision-making and preventing information loss between processing stages.

\subsection{Strategy Selection and Fallback Mechanisms}
\label{strategy-fallback}
A key innovation in our multi-agent architecture is the intelligent strategy selection and robust fallback mechanisms that enable the system to adapt to different query types and handle processing failures gracefully.

% \begin{figure}[htbp]
%     \centering
%     \includegraphics[width=0.9\linewidth]{assets/diagrams/strategy_fallback.png}
%     \caption{Strategy Selection and Fallback Mechanisms}
%     \label{fig:strategy-fallback}
% \end{figure}

As depicted in Figure~\ref{fig:multi-agent-architecture}, the strategy selection process begins with analyzing the question complexity and entity structure. The Strategy Selection Agent applies a decision-making algorithm to determine the optimal processing path:

\begin{enumerate}
    \item \textbf{Complexity Analysis}: The system analyzes question complexity by examining several factors:
    \begin{itemize}
        \item \textit{Entity Count}: Questions with multiple entities are routed to the SPARQL generation path.
        \item \textit{Relationship Indicators}: Phrases like ``and," ``or," ``as well as" suggest relationships between multiple entities, triggering SPARQL generation.
        \item \textit{User Settings}: Runtime settings can force specific processing paths based on user preferences.
    \end{itemize}
    
    \item \textbf{Adaptive Processing}: Based on complexity analysis, the system selects between:
    \begin{itemize}
        \item \textit{Verbalization Path}: For simple, single-entity questions where direct property lookup is likely to succeed.
        \item \textit{SPARQL Generation Path}: For complex questions involving multiple entities, relationships, or constraints.
    \end{itemize}
    
    \item \textbf{Fallback Mechanisms}: The system implements multiple fallback layers to ensure robust performance:
    \begin{itemize}
        \item \textit{Verbalization to SPARQL Fallback}: If verbalization fails to find relevant properties or returns insufficient results, the system automatically falls back to the SPARQL generation path.
        \item \textit{SPARQL to Google Search Fallback}: If SPARQL query generation or execution fails, and the Google Search option is enabled, the system falls back to web search to provide relevant information from external sources.
    \end{itemize}
\end{enumerate}

These fallback mechanisms significantly enhance the system's robustness by providing alternative processing paths when the primary approach fails. This multi-layered approach ensures that the system can deliver useful responses even in challenging scenarios where individual components might fail.

The strategy selection and fallback mechanisms are configurable through runtime settings, allowing users to control:

\begin{itemize}
    \item Whether to attempt verbalization before SPARQL generation
    \item Whether to enable Google Search as a fallback option
    \item Whether to enable translation for non-English queries
    \item Which knowledge source to use as the data source
\end{itemize}

This configurability provides users with fine-grained control over the system's behavior, enabling them to optimize the trade-off between processing efficiency and result comprehensiveness based on their specific use case requirements.

\subsection{Visualization and Transparency}
\label{vis-transparency}
A significant enhancement in our multi-agent architecture is the comprehensive visualization and transparency system that provides detailed insights into the reasoning process. This addresses the ``black box" nature of many AI systems by making the decision-making process explicit and reviewable.

The visualization system captures detailed logs of each agent's operations, including:

\begin{itemize}
    \item Input and output at each processing stage
    \item Decision points and routing logic
    \item Entity and property extraction results
    \item Generated SPARQL queries and their execution results
    \item Fallback triggers and alternative processing paths
    \item Timing information for performance analysis
\end{itemize}

These logs are stored in a semantic format using RDF (Resource Description Framework) triples, enabling sophisticated querying and analysis. The RDF-based logging system captures not only the sequence of events but also their semantic relationships, creating a knowledge graph of the processing workflow itself.

The visualization outputs include:

\begin{enumerate}
    \item \textbf{Real-time Reasoning Traces}: Detailed logs of the reasoning process as it unfolds, providing insights into how the system interprets and processes the question.
    
    \item \textbf{Process Flow Diagrams}: Visual representations of the processing workflow, showing how the question flows through different agents and the decisions made at each step.
    
    \item \textbf{RDF Knowledge Graphs}: Semantic representations of the processing workflow that can be queried and analyzed using SPARQL.
\end{enumerate}

This comprehensive visualization and transparency system enables users to understand and verify the system's reasoning process, enhancing trust and providing valuable insights for debugging and optimization.


\section{Evaluation Metrics}
\label{eval-metrics}

In this section, we present the comprehensive methodology used to evaluate our enhanced GraphRAG system. Our evaluation framework is designed to assess the performance of both our base model and the fine-tuned models across different knowledge domains and query complexity levels. We employ multiple complementary metrics to provide a holistic view of system performance.

\begin{figure}[htbp]
    \centering
    \includegraphics[width=0.8\linewidth]{assets/pics/evalframework.jpg}
    \caption{Comprehensive Evaluation Framework showing the components evaluated and metrics used at each stage}
    \label{fig:eval-framework}
\end{figure}

\subsection{SPARQL Generation Evaluation}
The evaluation of SPARQL query generation involves both syntactic and semantic assessment:

\begin{itemize}
    \item \textbf{Syntax Validity Rate}: The percentage of generated SPARQL queries that are syntactically valid and can be executed without errors
    \item \textbf{Average Generation Attempts}: The average number of generation attempts required to produce a valid query (with a maximum limit of 5 attempts per question)
\end{itemize}

For queries that require multiple generation attempts, we analyze the types of errors encountered and how the model adapts its generation strategy based on error feedback.

\subsection{Result Evaluation}
The most critical aspect of our evaluation is the assessment of query results. We use the following metrics to compare the results returned by generated queries against those from ground truth queries:

\begin{itemize}
    \item \textbf{Jaccard Similarity}: The size of the intersection divided by the size of the union of the result sets, measuring overall similarity
    \item \textbf{Precision}: The proportion of correct results returned by the generated query out of all results it returns
    \item \textbf{Recall}: The proportion of correct results returned by the generated query out of all results in the ground truth
    \item \textbf{F1 Score}: The harmonic mean of precision and recall, providing a balanced measure of result quality
\end{itemize}

These metrics are calculated using the following formulations:

Let $S_1$ be the set of results from the ground truth query and $S_2$ be the set of results from the generated query. Then:

\begin{align}
\text{Jaccard} &= \frac{|S_1 \cap S_2|}{|S_1 \cup S_2|} \\
\text{Precision} &= \frac{|S_1 \cap S_2|}{|S_2|} \\
\text{Recall} &= \frac{|S_1 \cap S_2|}{|S_1|} \\
\text{F1} &= \frac{2 \times \text{Precision} \times \text{Recall}}{\text{Precision} + \text{Recall}}
\end{align}

\subsection{Verbalization Evaluation}
For the verbalization approach, which provides direct natural language answers without generating SPARQL queries, we use:

\begin{itemize}
    \item \textbf{Verbalization Score}: A similarity score (0.0-1.0) that measures the semantic alignment between the question and the verbalized response
    \item \textbf{Verbalization Usage Rate}: The percentage of questions that were successfully answered using verbalization (verbalization score greater or equal than a certain threshold) without requiring SPARQL generation
\end{itemize}

\subsection{Stratified Evaluation}
To provide deeper insights into system performance across different scenarios, we stratify our evaluation based on the following approach:

\begin{itemize}
    \item \textbf{Knowledge Graph Domain}: We evaluate performance separately on each knowledge graph:
    \begin{itemize}
        \item Wikidata (using QALD-9-Plus without complexity stratification)
        \item Curriculum KG (with complexity stratification)
        \item Legal KG (with complexity stratification)
        \item GESIS KG (with complexity stratification)
    \end{itemize}
    
    \item \textbf{Query Complexity (Domain-Specific KGs Only)}: For the three domain-specific knowledge graphs (Curriculum KG, Legal KG, and GESIS KG), we categorize queries into three complexity levels:
    \begin{itemize}
        \item \textbf{Basic}: Simple one-hop queries retrieving direct properties of entities
        \item \textbf{Intermediate}: Two-hop queries or queries with simple constraints or aggregations
        \item \textbf{Advanced}: Multi-hop queries with complex constraints, unions, or aggregations
    \end{itemize}
\end{itemize}

For Wikidata evaluation, we do not apply query complexity stratification. While the QALD-9-Plus dataset naturally contains queries of varying complexity, we evaluate them as a single unified set without complexity-based categorization.

This stratification allows us to identify specific strengths and weaknesses of our system and to assess the effectiveness of our fine-tuning approach across different complexity levels (for domain-specific KGs) and knowledge domains.

\begin{figure}[htbp]
    \centering
    \includegraphics[width=0.6\linewidth]{assets/pics/compareval.jpg}
    \caption{Comparative Evaluation Methodology showing how we evaluate base and fine-tuned models across different dataset types and metrics}
    \label{fig:comparative-eval}
\end{figure}

\subsection{Evaluation Datasets}
For each knowledge graph domain, we employ dedicated evaluation datasets:

\begin{itemize}
    \item \textbf{Wikidata}: We utilize the QALD-9-Plus dataset, which contains 558 diverse natural language questions paired with ground truth SPARQL queries and answers. While we enhance this dataset with thoughts and entity/property matches for training purposes, evaluation is performed on the original question-SPARQL pairs without complexity stratification.
    
    \item \textbf{Curriculum KG}: We employ both template-based and random-walk generated questions covering basic course information, prerequisites, credits, and research groups, stratified across three complexity levels (basic, intermediate, advanced).
    
    \item \textbf{Legal KG}: We utilize both template-based and random-walk generated questions about legal document metadata, hierarchical relationships, amendments, and citations, organized by complexity level (basic, intermediate, advanced).
    
    \item \textbf{GESIS KG}: We employ both template-based and random-walk generated questions about scholarly publications, authors, research topics, and citation networks, stratified across three complexity levels (basic, intermediate, advanced).
\end{itemize}

All three domain-specific knowledge graphs (Curriculum KG, Legal KG, and GESIS KG) use a combination of template-based and random-walk generation methods to ensure comprehensive coverage of query patterns and linguistic diversity, with explicit complexity stratification for systematic evaluation.

\subsection{Evaluation Protocol}
Our evaluation protocol implements a systematic approach to assess all aspects of the GraphRAG system, with parallel evaluation tracks for both base and fine-tuned models:

\begin{enumerate}
    \item For each test question, we process it through two independent pipelines (base model pipeline and fine-tuned model pipeline).
    
    \item In each pipeline:
    \begin{enumerate}
        \item We extract entities and properties using the respective model (base or fine-tuned extraction model).
        \item We attempt verbalization for each identified entity and record verbalization scores.
        \item Based on the verbalization score for that specific pipeline, we either:
        \begin{itemize}
            \item Accept the verbalization result if the score exceeds the threshold, or
            \item Generate a SPARQL query using the respective SPARQL generation model (base or fine-tuned) if the verbalization score falls below the threshold.
        \end{itemize}
        \item We execute the resulting query (either from verbalization or SPARQL generation) against the appropriate knowledge graph and collect results.
    \end{enumerate}
    
    \item We calculate all evaluation metrics by comparing each pipeline's results against ground truth results.
    
    \item We stratify results by knowledge graph domain for all KGs. For domain-specific knowledge graphs (Curriculum KG, Legal KG, and GESIS KG), we additionally stratify by query complexity (basic, intermediate, advanced). For Wikidata, we only report domain-level results without complexity breakdown.
\end{enumerate}

This parallel evaluation approach allows us to directly compare the performance of base and fine-tuned models, even when they take different processing paths (verbalization vs. SPARQL generation) for the same question.

For reproducibility, we document the complete evaluation results, including all generated queries, extraction outputs, verbalization scores, and execution results, along with detailed error analyses for failed queries.

Throughout the model adaptation process, we also monitor multiple metrics to assess model performance:

\begin{itemize}
    \item \textbf{Learning curves}: Training and validation loss progression to detect overfitting and ensure proper convergence
    \item \textbf{Query execution reliability}: Percentage of generated queries that execute without errors
    \item \textbf{Answer accuracy}: Comparison of query results with ground truth answers
\end{itemize}

Models are evaluated on held-out validation sets, with the best-performing model (based on validation metrics) selected for comprehensive evaluation. This methodical evaluation approach ensures that the fine-tuned models maintain high quality across different domains and query types.

The combined metrics from both component-level and end-to-end evaluation provide a comprehensive assessment of our enhanced GraphRAG system, enabling detailed comparison with the original baseline and quantitative validation of our methodological improvements across different knowledge graph domains and query complexity levels.